Tables~\ref{tab:graphtypes:one}-\ref{tab:graphtypes:two} 
list Orbiter commands to create graphs.
\orbitertableofcommands{create_graph_1}
\orbitertableofcommands{create_graph_2}



\bigskip

Let us consider some examples.

\bigskip 

The command 

\sourcecode{make_graph_from_adjacency_matrix}

creates a graph on three vertices from a given adjacency matrix.
The adjacency matrix must be given as a vector object, with the matrix entries listed in row-major ordering.
In the example, the adjacency matrix is 
$$
\left[
\begin{array}{cccc}
0&1&1 \\
1&0&1 \\
1&1&0
\end{array}
\right]
$$
which is given as a vector 
$$
0,1,1, 1,0,1, 1,1,0.
$$

\bigskip 


The adjacency matrix of a graph can be read from a file. 
The following command sequence 
uses a makefile variable to store the adjacency matrix of a graph. 
The matrix is then copied into a file and the graph is created from the file. 
Here is the makefile variable containing the adjacency matrix:

\sourcecodevariable{TRIANGLE_GRAPH}

The following command writes the adjacency matrix 
to a csv file and then creates 
the graph from the csv file:

\sourcecode{make_triangle_graph}

This will create the same triangle graph that we have created before. 



\bigskip


The command 

\sourcecode{Cycle_graph_13}

creates the cycle graph of degree 13. 

\bigskip 



The command 

\sourcecode{Chain_232}

creates the chain graph with respect to the partitions $(2,3,2)$ and $(2,3,2)$. 



\bigskip

The command 

\sourcecode{Paley_13_graph}

creates the Paley graph on 13 vertices. 
The command first creates the field $\bbF_{13}$ 
and then the associated Payley graph.

\bigskip

The command 

\sourcecode{Paley_13_graph_complement}

creates the complement of the Paley graph on 13 vertices. 

\bigskip

The command 

\sourcecode{Winnie_Li_graph_q16_index2}

creates the Winnie Li graph~\cite{Li1992} 
on 16 vertices with index 2.
The command first creates the field $\bbF_{16}$ 
and then the associated graph.

\bigskip


The command 

\sourcecode{Sarnak_graph_29_5}

creates the graph $X_{29,5}$ from~\cite{Lubotzky1988} on 60 vertices.

\bigskip



The command 

\sourcecode{Neumaier_graph_16}

creates a Neumaier graph on 16 vertices.

\bigskip

The command 

\sourcecode{tritangent_planes_graph}

creates the graph of tritangent planes in a general cubic surface with 27 lines. 
Two vertices are adjacent if the associated tritangent planes are line-disjoint.


\bigskip


The command 

\sourcecode{trihedral_pair_graph}

creates the graph of trihedral pairs in a general cubic surface with 27 lines. 
Two vertices are adjacent if the associated trihedral pairs are line-disjoint.


\bigskip


The command 

\sourcecode{small_graph}

creates a graph by listing the edges in pairs. 
The following graph is created:
\orbiteroutput{GRAPHICS/WRAPPERS/a_small_graph}



\bigskip

The command 

\sourcecode{petersen}

creates the Petersen graph. 

\bigskip

The command 

\sourcecode{Johnson_6_2_0}

creates the Johnson graph $J(6,2,0).$ 

\bigskip


The command 

\sourcecode{Hamming_graph_3}

creates the Hamming graph of order 3. 

\bigskip

The command 

\sourcecode{Hamming_graph_7}

creates the Hamming graph of order 7. 

\bigskip


The command 

\sourcecode{Op_6_2_coll_graph}

creates the collinearity graph of the orthogonal space
$O^+(6,2).$






\bigskip

There is a graph on 315 vertices 
that arises from the Cohen-Tits near octagon (see~\cite{Brouwer315}). 
The graph was first constructed in~\cite{Cohen1981} 
and has automorphism group equal to $\Aut(HJ)$, 
the automorphism group of the Higman-Sims sporadic simple group.
The graph is distance-regular.
An incidence matrix can be found in 
Ascii format on the web site~\cite{Bayley}. 
In the following, we assume that a file \verb'halljanko315.csv' 
is present, containing the incidence matrix of the graph. 
The following command creates the graph from the file:

\sourcecode{HJ_graph}

In Section~\ref{sec:canonical:form:objects:graphs}, 
we will compute the automorphism group of the graph (of order 1209600). 
This will create a file \verb'halljanko315_gens.csv' 
which we use here to create an orbital graph.
An orbital graph is a graph whose adjacency matrix corresponds 
to an orbit of a permutation group in the action on pairs. 
The group is the automorphism group of the graph.
The following command creates the third orbital graph:

\sourcecode{Group_Perm315_Orbital_3.colored_graph}


\bigskip

Table~\ref{tab:graph:modifiers} 
shows some Orbiter commands to modify graphs.
The commands replace the given graph 
by the graph obtained after applying the specified modification.
\orbitertableofcommands{graph_modification}

\bigskip


For a graph $\Gamma$, the distance 2 graph $\Delta$ 
has the same vertices as $\Gamma$, 
and two vertices in $\Delta$ are adjacent 
if and only if the distance in $\Gamma$ is two.
The following command creates 
the distance 2 graph of the Cohen-Tits graph. 

\sourcecode{HJ_d2_graph}




\bigskip


Let us look at some examples of Cayley graphs. 
The first graph has $G = \bbZ_{11}$ 
and connection set all elements congruent $1$ mod $3$. 
We create the group as a subgroup of 
the one-dimensional affine group over $\bbF_{11}.$ 
This means that the map $x \mapsto ax+b$ mod $11$ 
is encoded as a pair $(a,b).$  

\sourcecode{Cayley_Z11_1mod3}

The vertices of the Cayley graph are ordered. 
The ordering is determined by the stabilizer chain. 
This is a total ordering.

\bigskip


In the following example, we create a Cayley graph 
based on the symmetric group on 4 things. 
We take the Coxeter generators as connection set:

\sourcecode{Cayley_Sym4_coxeter}

The star graph is another Cayley graph for the symmetric group. 
The connection set is given by the permutations 
$(0,\,i)$ for $i=1,\ldots,n-1.$
The next example creates the star graph on 4 vertices:

\sourcecode{Cayley_Sym4_star}



\bigskip

In the following example, 
we create a graph of nonattacking queens on a chess board.

\sourcecode{Queens_graph}

\bigskip


In the following example, 
we create an affine polar graph based on the $O^+(4,2)$ space.

\sourcecode{Op_4_2_affine_polar_graph}

\bigskip


Next, we create the first 
and then the 15th graph in the list of Paulus and Rozenfeld.

\sourcecode{Paulus_1}

\bigskip

\sourcecode{Paulus_15}

\bigskip

The next command creates the Higman Sims graph. The graph is on 100 points.
We begin with the Mathieu group $M_{24}$, which arises as the stabilizer of the Golay code.
The Golay code is a projective code of length 24. 
The colums of the generator matrix correspond to a set $S$ of 
24 points in $\PG(11,2)$ with the property that any 7 are independent. 
This set of points can be encoded using Orbiter point ranks as the following set:

\sourcecodevariable{GOLAY_24_CODE_PROJECTIVE}

Next, we need precomputed generator for the group $M_{24}$. These are a set of matrices of size $12 \times 12$ acting on $\PG(11,2)$ and stabilizing the set $S$. They
are encoded as a makefile variable: 

\sourcecodevariable{GOLAY_24_CODE_AUT_GENS}

The following code snippet creates the graph:

\sourcecode{Higman_Sims_graph}


In lines 7-9, we create the set $S$ of size 24. In line 14, we create the field $\bbF_2$. 
In lines 15-18, we create the group $M_{24}$ as a matrix subgroup of $\PGL(12,2).$
In lines 19-21, we create the permutation representation of $M_{24}$ of degree 24 
acting on the set $S$.
In lines 22-24, we create the subgroup $M_{22}.2$ 
which stabilizes the elements $\{0,1\}$ of $S$ as a set. 
In lines 25-28, we create the permutation action of $M_{22}.2$ 
on the 22 points of $S$ 
different from $\{0,1\}$. 
Note that the action is re-indexing the point set as $0-21.$
In lines 37-42, we invoke the poset classification algorithm 
to compute the orbits of $M_{22}.2$ 
on subsets of size at most 6. 
This computation takes place in the action of degree 22 just constructed. 
In lines 43-45, we create the Kramer-Mesner matrix for the orbits on $3$-subsets versuc the orbits on $6$-subsets. 
This matrix allows us to identify the orbit on $6$-subsets which forms 
a Steiner 3-design. 
In the numbering chosen by the poset classification algorithm, 
it is the second orbit. 
Since indexing is zero-based, this is orbit 1.
In lines 55-57, we create the design from this orbit.
This is the Witt design on 22 points. 
In lines 58-61, we export the design to file by means of the set of flags. 
In lines 62-65, we compute the intersection matrix of the design. 
This matrix can be used to give an independent verification that the chosen orbit has the design property. 
In lines 70-73, we compute the orbits of the group of the design (which is $M_{11}.2$) 
on pairs of blocks. 
In lines 74-77, we create the one-point extension of the Witt design based on the orbit on pairs of index 1. 
Finally, in lines 78-81, the Higman-Sims graph is constructed from the one-point extension. 










\bigskip

